\section*{ID9509: Geometry: Curvature of Spheres and Tori: (R)/(R) ≈ 2.46}
\paragraph{Metadata.}
PiID: 1908699603302 | RTEID: RTE:P36925.5587:T4.2933 | UUID: 2f177b50-24b0-514d-ae9f-fddd9e0ff091

\paragraph{Canonical Statement.}
\$) one finds \$K\_\{\textbackslash{}text\{outer\}\} = \textbackslash{}frac\{1\}\{R, (R+r)\}\$, while at the inner equator (\$\textbackslash{}phi=\textbackslash{}pi\$) \$K\_\{\textbackslash{}text\{inner\}\} = -,\textbackslash{}frac\{1\}\{R, (R-r)\}\$. The \textbackslash{}emph\{intrinsic\} curvature thus differs between outer and inner regions, even though the torus is embedded in flat 3D space (its \textbackslash{}emph\{extrinsic\} curvature is encoded in how it bends in space). The ratio of magnitudes of inner to outer curvature is \textbackslash{}begin\{equation\} \textbackslash{}label\{eq:torus-curv-ratio\} \textbackslash{}frac\{|K\_\{\textbackslash{}text\{inner\}\}|\}\{K\_\{\textbackslash{}text\{outer\}\}\} ;=; \textbackslash{}frac\{R + r\}\{,R - r,\},, \textbackslash{}end\{equation\} a dimensionless number purely determined by geometry. Notably, if one chooses \$\textbackslash{}frac\{R\}\{r\} \textbackslash{}approx 2.46\$, this curvature ratio evaluates to \$\textbackslash{}approx 2.37037\$ – precisely the multiple \$2\textbackslash{}chi\$. In other words, a torus whose major radius is about \$2.46\$ times its tube radius will exhibit an inner/outer curvature contrast of 237.037\%. This is a concrete geometrical realization of \$\textbackslash{}chi\$’s appearance: the factor by which intrinsic curvature “stretches” between different regions of a curved surface.

\paragraph{Canonical Equation.}
\[
\frac{R}{r} \approx 2.46
\]

\paragraph{Dictionary.}
Geometry: Curvature of Spheres and Tori: (R)/(R) ≈ 2.46 — (R)/(r) ≈ 2.46

\paragraph{Encyclopedia.}
Geometry: Curvature of Spheres and Tori: (R)/(R) ≈ 2.46 is a scaling / order-of-magnitude relation, written (R)/(r) ≈ 2.46. It expresses a scaling or proportionality, capturing how the quantity grows with its parameters. It is built from R, r. Within the Trawin Zero Point registry it serves as one element of the framework's mathematical narrative.
